\begin{tomtat}
	\begin{minipage}[t]{0.3\textwidth}
		\vspace{0pt}
		\begin{tomtat}
			\begin{center}
				\textbf{ĐỀ 05}
			\end{center}
			
		\end{tomtat}
	\end{minipage}
	\hfill
	\begin{minipage}[t]{0.5\textwidth}
		\vspace{0pt}
		\begin{tomtat}
			\begin{center}
				\includegraphics[scale=0.08]{images/5}
			\end{center}
		\end{tomtat}
	\end{minipage}
\end{tomtat}
\textbf{I. PHẦN TRẮC NGHIỆM}
\begin{ex}
	Rút gọn biểu thức $A=3\sqrt{20}-\sqrt{500}+\sqrt{5}$. ta được
	\choice{$-3\sqrt{5}$}{$\sqrt{5}$}{$3\sqrt{5}$}{$2\sqrt{5}$}
	\loigiai{
		Ta có:
		\begin{align*}
			A &= 3\sqrt{20}-\sqrt{500}+\sqrt{5} \\
			&= 3\sqrt{4 \cdot 5} - \sqrt{100 \cdot 5} + \sqrt{5} \\
			&= 3 \cdot 2\sqrt{5} - 10\sqrt{5} + \sqrt{5} \\
			&= 6\sqrt{5} - 10\sqrt{5} + \sqrt{5} \\
			&= (6 - 10 + 1)\sqrt{5} \\
			&= -3\sqrt{5}
		\end{align*}
		\\ \textbf{ĐÁP ÁN A}
	}
\end{ex}

\begin{ex}
	Cho tam giác $ABC$ vuông tại $A$. Biết độ dài cạnh $AB = 6 \text{ cm}, AC = 8 \text{ cm}$. Thì $\sin \widehat{ABC}$ bằng
	\choice{0,75}{0,8}{0,6}{0,7}
	\loigiai{
		Trong tam giác $ABC$ vuông tại $A$, theo định lý Pitago, ta có:
		$BC^2 = AB^2 + AC^2$
		$BC^2 = 6^2 + 8^2$
		$BC^2 = 36 + 64$
		$BC^2 = 100$
		$BC = \sqrt{100} = 10 \text{ cm}$.
		
		Theo định nghĩa tỉ số lượng giác sin trong tam giác vuông, ta có:
		$$ \sin \widehat{ABC} = \sin B = \frac{\text{cạnh đối}}{\text{cạnh huyền}} = \frac{AC}{BC} $$
		Với $AC=8 \text{ cm}$ và $BC=10 \text{ cm}$, ta có:
		$$ \sin \widehat{ABC} = \frac{8}{10} = 0,8 $$
		\\ \textbf{ĐÁP ÁN B}
	}
\end{ex}

\begin{ex}
	Tính thể tích của một viên bi hình cầu có bán kính $6 \text{ cm}$ (lấy $\pi \approx 3,14$ và làm tròn kết quả đến hàng phần mười).
	\choice{$903 \text{ cm}^3$}{$904 \text{ cm}^3$}{$804 \text{ cm}^3$}{$905 \text{ cm}^3$}
	\loigiai{
		Thể tích của hình cầu được tính bằng công thức: $V = \frac{4}{3}\pi R^3$.
		Trong đó $R$ là bán kính của hình cầu.
		Theo đề bài, $R = 6 \text{ cm}$ và $\pi \approx 3,14$.
		\begin{align*}
			V &= \frac{4}{3} \times 3,14 \times (6)^3 \\
			&= \frac{4}{3} \times 3,14 \times 216 \\
			&= 4 \times 3,14 \times \frac{216}{3} \\
			&= 4 \times 3,14 \times 72 \\
			&= 12,56 \times 72 \\
			&= 904,32 \text{ cm}^3
		\end{align*}
		Làm tròn kết quả đến hàng phần mười, ta được $904,3 \text{ cm}^3$.
		(Nếu đề yêu cầu làm tròn đến hàng đơn vị thì gần với đáp án B nhất)
		Kiểm tra lại các lựa chọn, có vẻ đề bài muốn làm tròn đến hàng đơn vị.
		$904,32 \text{ cm}^3$ làm tròn đến hàng đơn vị là $904 \text{ cm}^3$.
		\\ \textbf{ĐÁP ÁN B}
	}
\end{ex}

\begin{ex}
	Diện tích xung quanh của hình trụ có bán kính đáy $20 \text{ cm}$ và chiều cao $30 \text{ cm}$ (lấy $\pi \approx 3,14$ và làm tròn kết quả đến hàng phần mười) là
	\choice{$3768 \text{ cm}^2$}{$3769 \text{ cm}^2$}{$2768 \text{ cm}^2$}{$3768 \text{ cm}^2$}
	\loigiai{
		Diện tích xung quanh của hình trụ được tính bằng công thức: $S_{xq} = 2\pi Rh$.
		Trong đó:
		$R$ là bán kính đáy ($R=20 \text{ cm}$).
		$h$ là chiều cao ($h=30 \text{ cm}$).
		Sử dụng $\pi \approx 3,14$.
		\begin{align*}
			S_{xq} &= 2 \times 3,14 \times 20 \times 30 \\
			&= 6,28 \times 600 \\
			&= 3768 \text{ cm}^2
		\end{align*}
		Kết quả là $3768 \text{ cm}^2$, không cần làm tròn thêm vì đã là số nguyên.
		\\ \textbf{ĐÁP ÁN A} (Lưu ý: Đáp án D giống A)
	}
\end{ex}

\begin{ex}
	Tìm tất cả các nghiệm của bất phương trình $6 - 2x > 0$.
	\choice{$x > 3$}{$x > -3$}{$x < 3$}{$x < -3$}
	\loigiai{
		Ta giải bất phương trình:
		\begin{align*}
			6 - 2x &> 0 \\
			-2x &> -6
		\end{align*}
		Khi chia cả hai vế của bất phương trình cho một số âm (ở đây là -2), ta phải đổi chiều bất đẳng thức:
		\begin{align*}
			x &< \frac{-6}{-2} \\
			x &< 3
		\end{align*}
		Vậy nghiệm của bất phương trình là $x < 3$.
		\\ \textbf{ĐÁP ÁN C}
	}
\end{ex}

\begin{ex}
	Một vật rơi tự do từ độ cao so với mặt đất là $120 \text{ m}$. Bỏ qua sức cản không khí, quãng đường chuyển động $s$ (mét) của vật rơi sau thời gian $t$ được biểu diễn gần đúng bởi công thức $s = 5t^2$, trong đó $t$ là thời gian tính bằng giây. Hỏi sau 4 giây vật này cách mặt đất bao nhiêu mét ?
	\choice{$40 \text{ m}$}{$80 \text{ m}$}{$100 \text{ m}$}{$90 \text{ m}$}
	\loigiai{
		Quãng đường vật rơi được sau thời gian $t=4$ giây là:
		$$ s = 5t^2 = 5 \times (4)^2 = 5 \times 16 = 80 \text{ m} $$
		Đây là quãng đường vật đã rơi từ vị trí ban đầu.
		Độ cao ban đầu của vật so với mặt đất là $120 \text{ m}$.
		Vậy, sau 4 giây, vật này cách mặt đất một khoảng là:
		$$ \text{Khoảng cách} = \text{Độ cao ban đầu} - \text{Quãng đường rơi} $$
		$$ \text{Khoảng cách} = 120 \text{ m} - 80 \text{ m} = 40 \text{ m} $$
		\\ \textbf{ĐÁP ÁN A}
	}
\end{ex}

\begin{ex}
	Điểm bài kiểm tra môn toán của 30 học sinh lớp 9A được Thầy giáo ghi lại như sau:
	\begin{center}
		\begin{tabular}{|c|c|c|c|c|c|c|c|c|c|}
			\hline
			9 & 8 & 10 & 7 & 8 & 9 & 7 & 10 & 6 & 8 \\ \hline
			8 & 7 & 7 & 10 & 8 & 10 & 7 & 6 & 8 & 7 \\ \hline
			10 & 8 & 10 & 9 & 6 & 8 & 8 & 7 & 6 & 9 \\ \hline
		\end{tabular}
	\end{center}
	Tần số của giá trị 8 điểm là
	\choice{7}{8}{100}{9}
	\loigiai{
		Để tìm tần số của giá trị 8 điểm, ta đếm số lần xuất hiện của điểm 8 trong bảng:
		Dòng 1: 8 (1 lần), 8 (2 lần), 8 (3 lần), 8 (4 lần)
		Dòng 2: 8 (5 lần), 8 (6 lần), 8 (7 lần), 8 (8 lần), 8 (9 lần)
		Vậy, điểm 8 xuất hiện 9 lần.
		Tần số của giá trị 8 điểm là 9.
		\\ \textbf{ĐÁP ÁN D}
	}
\end{ex}

\begin{ex}
	Người ta thống kê 4 loại ô tô chạy qua một trạm thu phí trong một giờ và vẽ được biểu đồ tần số như hình bên dưới
	\begin{center}
		\begin{tabular}{|l|c|c|c|c|}
			\hline
			Loại xe & 4 chỗ & 7 chỗ & 9 chỗ & 16 chỗ \\
			\hline
			Số xe   & 10    & 17    & 8     & 5      \\
			\hline
		\end{tabular}
	\end{center}
	Từ bảng biểu đồ tần số, hãy cho biết loại xe nào đi qua trạm thu phí nhiều nhất.
	\choice{7 chỗ}{7 chỗ}{9 chỗ}{16 chỗ}
	\loigiai{
		Dựa vào bảng tần số:
		\begin{itemize}
			\item Xe 4 chỗ: 10 xe
			\item Xe 7 chỗ: 17 xe
			\item Xe 9 chỗ: 8 xe
			\item Xe 16 chỗ: 5 xe
		\end{itemize}
		So sánh số lượng xe của các loại, ta thấy loại xe 7 chỗ có số lượng nhiều nhất (17 xe).
		\\ \textbf{ĐÁP ÁN A} (Lưu ý: Đáp án B giống A)
	}
\end{ex}

\textbf{II. PHẦN TỰ LUẬN}
\begin{ex}
	Cho $x \ge 0, x \ne 1$. Rút gọn biểu thức $P = \left(\frac{7-\sqrt{x}}{x-1} + \frac{3}{\sqrt{x}+1}\right) \cdot \frac{\sqrt{x}+2}{\sqrt{x}-1}$
	\loigiai{
		Với $x \ge 0, x \ne 1$, ta có:
		\begin{align*}
			P &= \left(\frac{7-\sqrt{x}}{(\sqrt{x}-1)(\sqrt{x}+1)} + \frac{3(\sqrt{x}-1)}{(\sqrt{x}+1)(\sqrt{x}-1)}\right) \cdot \frac{\sqrt{x}+2}{\sqrt{x}-1} \\
			&= \frac{7-\sqrt{x} + 3\sqrt{x}-3}{(\sqrt{x}-1)(\sqrt{x}+1)} \cdot \frac{\sqrt{x}+2}{\sqrt{x}-1} \\
			&= \frac{2\sqrt{x}+4}{(\sqrt{x}-1)(\sqrt{x}+1)} \cdot \frac{\sqrt{x}+2}{\sqrt{x}-1} \\
			&= \frac{2(\sqrt{x}+2)}{(\sqrt{x}-1)(\sqrt{x}+1)} \cdot \frac{\sqrt{x}+2}{\sqrt{x}-1} \\
			&= \frac{2(\sqrt{x}+2)^2}{(\sqrt{x}-1)^2(\sqrt{x}+1)}
		\end{align*}
		Vậy $P = \frac{2(\sqrt{x}+2)^2}{(\sqrt{x}-1)^2(\sqrt{x}+1)}$.
	}
\end{ex}

\begin{ex}
	Giải hệ phương trình 	$ \begin{cases} 2x - y = 2 \\ x + y = 7 \end{cases} $
	\loigiai{
		Ta có hệ phương trình:
		$$ \begin{cases} 2x - y = 2 & (1) \\ x + y = 7 & (2) \end{cases} $$
		Cộng vế theo vế của phương trình (1) và phương trình (2), ta được:
		\begin{align*}
			(2x - y) + (x + y) &= 2 + 7 \\
			3x &= 9 \\
			x &= 3
		\end{align*}
		Thay $x=3$ vào phương trình (2):
		\begin{align*}
			3 + y &= 7 \\
			y &= 7 - 3 \\
			y &= 4
		\end{align*}
		Vậy nghiệm của hệ phương trình là $(x; y) = (3; 4)$.
	}
\end{ex}

\begin{ex}
	Phương trình $x^2 - 4x - 1 = 0$ có hai nghiệm phân biệt $x_1, x_2$. Tính giá trị của $T = x_1^2 + 4x_2 - x_1x_2$.
	\loigiai{
		Phương trình đã cho là $x^2 - 4x - 1 = 0$.
		Đây là phương trình bậc hai có dạng $ax^2+bx+c=0$ với $a=1$, $b=-4$, $c=-1$.
		Ta tính biệt thức $\Delta = b^2-4ac = (-4)^2 - 4(1)(-1) = 16 + 4 = 20$.
		Vì $\Delta = 20 > 0$ nên phương trình có hai nghiệm phân biệt $x_1, x_2$.
		
		Theo hệ thức Vi-ét, ta có:
		$$ \begin{cases} 
			x_1+x_2 = \dfrac{-b}{a} = \dfrac{-(-4)}{1} = 4 \\ 
			x_1x_2 = \dfrac{c}{a} = \dfrac{-1}{1} = -1 
		\end{cases} $$
		Vì $x_1$ là nghiệm của phương trình $x^2 - 4x - 1 = 0$, nên $x_1^2 - 4x_1 - 1 = 0 \Rightarrow x_1^2 = 4x_1 + 1$.
		Thay vào biểu thức $T$:
		\begin{align*}
			T &= x_1^2 + 4x_2 - x_1x_2 \\
			&= (4x_1 + 1) + 4x_2 - x_1x_2 \\
			&= 4(x_1+x_2) + 1 - x_1x_2
		\end{align*}
		Thay giá trị $x_1+x_2 = 4$ và $x_1x_2 = -1$ vào biểu thức $T$:
		\begin{align*}
			T &= 4(4) + 1 - (-1) \\
			&= 16 + 1 + 1 \\
			&= 18
		\end{align*}
		Vậy giá trị của biểu thức $T$ là 18.
	}
\end{ex}

\begin{ex}
	Hai người đi xe đạp cùng xuất phát một lúc ở cùng một địa điểm. Người thứ nhất đi về phía Tây, người thứ 2 đi về phía Bắc. Sau 2h họ cách nhau 60km theo đường chim bay. Tính tốc độ của mỗi người, biết tốc độ của người thứ 2 lớn hơn tốc độ của người thứ 1 là 6km/h.
	\loigiai{
		Gọi tốc độ của người thứ nhất là $x$ (km/h) và tốc độ của người thứ hai là $y$ (km/h).
		Điều kiện: $x > 0, y > 0$.
		Theo đề bài, tốc độ của người thứ 2 lớn hơn tốc độ của người thứ 1 là 6km/h, nên ta có phương trình:
		$$ y - x = 6 \quad \Rightarrow y = x + 6 $$
		Sau 2 giờ:
		Quãng đường người thứ nhất đi được là $S_1 = x \cdot 2 = 2x$ (km).
		Quãng đường người thứ hai đi được là $S_2 = y \cdot 2 = 2y = 2(x+6)$ (km).
		Vì người thứ nhất đi về phía Tây và người thứ hai đi về phía Bắc, nên quãng đường họ đi được tạo thành hai cạnh góc vuông của một tam giác vuông. Khoảng cách giữa họ sau 2 giờ là cạnh huyền của tam giác vuông đó.
		Theo định lý Pitago, ta có: $S_1^2 + S_2^2 = (\text{khoảng cách})^2$.
		$$ (2x)^2 + (2y)^2 = 60^2 $$
		$$ 4x^2 + 4y^2 = 3600 $$
		Chia cả hai vế cho 4:
		$$ x^2 + y^2 = 900 $$
		Ta có hệ phương trình:
		$$ \begin{cases} y = x + 6 \\ x^2 + y^2 = 900 \end{cases} $$
		Thay $y = x + 6$ vào phương trình thứ hai:
		\begin{align*}
			x^2 + (x+6)^2 &= 900 \\
			x^2 + x^2 + 12x + 36 &= 900 \\
			2x^2 + 12x + 36 - 900 &= 0 \\
			2x^2 + 12x - 864 &= 0
		\end{align*}
		Chia cả hai vế cho 2:
		$$ x^2 + 6x - 432 = 0 $$
		Tính biệt thức $\Delta = b^2 - 4ac = 6^2 - 4(1)(-432) = 36 + 1728 = 1764$.
		$\sqrt{\Delta} = \sqrt{1764} = 42$.
		Phương trình có hai nghiệm:
		$$ x_1 = \frac{-6 + 42}{2(1)} = \frac{36}{2} = 18 $$
		$$ x_2 = \frac{-6 - 42}{2(1)} = \frac{-48}{2} = -24 $$
		Vì $x > 0$, nên ta nhận $x = 18$ km/h.
		Khi đó, tốc độ của người thứ hai là $y = x + 6 = 18 + 6 = 24$ km/h.
		
		Vậy tốc độ của người thứ nhất là 18 km/h và tốc độ của người thứ hai là 24 km/h.
	}
\end{ex}

\begin{ex}
	Một hộp đựng 4 thẻ, mỗi thẻ được ghi một trong các số 1, 2, 3, 4; hai thẻ khác nhau thì ghi hai số khác nhau. Rút ngẫu nhiên lần lượt 2 thẻ từ hộp đó (mỗi lần rút thẻ ra không bỏ thẻ lại hộp). Tính xác suất của biến cố A: "Rút được 2 thẻ có tổng hai số ghi trên thẻ là một số chia hết cho 3".
	\loigiai{
		Tập hợp tất cả các kết quả có thể xảy ra là:\\
		$\Omega = \{$
		(1,2), (1,3), (1,4), 
		(2,1), (2,3), (2,4), 
		(3,1), (3,2), (3,4), 
		(4,1), (4,2), (4,3) 
		$\}$. \\
	Tập hợp các kết quả thuận lợi cho biến cố A là: \\ $A = \{(1,2), (2,1), (2,4), (4,2)\}$.
		Số phần tử của biến cố A là $n(A) = 4$.
		
		Xác suất của biến cố A là:
		$$ P(A) = \frac{n(A)}{n(\Omega)} = \frac{4}{12} = \frac{1}{3} $$
		Vậy xác suất của biến cố A là $\frac{1}{3}$.
	}
\end{ex}

\begin{ex}
	Cho đường tròn $(O)$ và dây cung $AB$ không đi qua tâm đường tròn. Gọi $H$ là trung điểm của $AB$, đường thẳng $HO$ cắt đường tròn $(O)$ tại hai điểm $C$ và $D$ ($C$ thuộc cung nhỏ $AB$). Điểm $M$ thuộc đoạn $AH$ ($M$ khác $A$ và $H$), tia $CM$ cắt đường tròn $(O)$ tại điểm thứ hai là $N$.
	\begin{itemize}
		\item[a)] Chứng minh 4 điểm $D, N, M, H$ cùng thuộc một đường tròn.
		\item[b)] Gọi $I$ là giao điểm của hai đường thẳng $DN$ và $AB$, $K$ là giao điểm của hai đường thẳng $CI$ và $DM$. Chứng minh rằng $CM \cdot CN = CK \cdot CI$ và $M$ là tâm đường tròn nội tiếp tam giác $HKN$.
	\end{itemize}
	\loigiai{
	}
		
\end{ex}
\setcounterref{ex}{0}